\documentclass[10pt,a4paper]{amsart}
\usepackage[margin=19mm]{geometry}
\usepackage{amsmath,amssymb,mathtools}
\usepackage[scr=rsfs,cal=euler]{mathalfa}
\usepackage{xfrac}
\usepackage[unicode,hidelinks,bookmarks=false]{hyperref}
\newcommand{\ie}{i.e.\ }
\newcommand{\cf}{cf.\ }
\newcommand{\resp}{respectively\ }
\newcommand{\Subsection}{\subsection}
\providecommand{\nobreakdash}{\nobreak\hbox{-}}
\setlength{\emergencystretch}{2em}
\multlinegap0pt
\usepackage{bm}
\usepackage{bbm}
\usepackage{dsfont}
\usepackage{xspace}
\usepackage{cancel}
\usepackage{mathtools}
\usepackage{cleveref}
\usepackage[shortlabels]{enumitem}
\usepackage{amsthm}
\makeatletter
\@ifundefined{theorem}{\newtheorem{theorem}{Theorem}}{}
\@ifundefined{assumption}{\newtheorem{assumption}[theorem]{Assumption}}{}
\makeatother
\newcommand{\E}{\mathbb{E}}
\newcommand{\Pcal}{\mathcal{P}}
\newcommand{\R}{\mathbb{R}}
\newcommand{\abs}[1]{\left\lvert#1\right\rvert}
\newcommand{\dd}{\mathrm{d}}
\newcommand{\norm}[1]{\left\lVert#1\right\rVert}
\newcommand{\scpr}[2]{\left\langle #1, #2 \right\rangle}
\title[Neural Brownian Motion]{Neural Brownian Motion}
\author{Qian Qi}
\date{Source: arXiv:2507.14499v1 (CC BY 4.0)}
\begin{document}
\maketitle

\noindent\emph{Source and licence.} Neural Brownian Motion. Version arXiv:2507.14499v1 (CC BY 4.0), distributed under the Creative Commons Attribution 4.0 International licence, as recorded in the arXiv licence metadata for this exact version and shown on the abstract page https://arxiv.org/abs/2507.14499v1. Full rights notice: licenses/arxiv-2507.14499-CC-BY-4.0.txt.\par\smallskip

\noindent\textit{Editorial scope.} Counted unit: the Theorem whose source label is \texttt{thm:mckean\_driver\_existence} in the article (source0.tex lines 690--693, source label \texttt{thm:mckean\_driver\_existence}), delivered together with the article's own complete original proof of it (source0.tex lines 695--782, one \texttt{proof} environment of 88 lines). No mathematics is written, repaired or re-derived: statement and proof are the article's own text line for line. Required context retained uncounted: ass:wellposedness\_multidim (lines 199--211); thm:existence\_of\_theta (lines 240--243); ass:implicit\_vol (lines 413--421); ass:mckean\_lipschitz\_revised (lines 671--683). Every definition, display and label of those lines is retained unchanged. Omitted from this package and not redistributed: 1--198, 212--239, 244--412, 422--670, 684--689, 694--694, 783--1168. Nothing inside a retained excerpt is omitted or altered: every retained body is delivered exactly as the article's lines are written, so no omission receipt of any kind accompanies this package. Citations: the retained excerpts carry no citation command at all, so no bibliography entry of the article is transcribed and no reference is redistributed. Reference adaptation: none was needed and none was made; every reference inside the delivered unit resolves inside it, so no unresolved reference or citation is produced. Printed slips kept literal: no printed word, symbol, hypothesis or number of the counted statement or of its proof was altered. Layout and class: the article's own class is \texttt{imsart} with packages including amssymb, amsfonts, lmodern, microtype, booktabs, enumitem, bm, xcolor, natbib, mathtools, graphicx, caption, subcaption, tabularx; the wrapper is the lane's pinned \texttt{amsart} editorial wrapper at 10pt on A4, which restates the theorem-like environments the retained text needs and adds the article's own macro definitions that the retained text needs (\texttt{\textbackslash E}, \texttt{\textbackslash Pcal}, \texttt{\textbackslash R}, \texttt{\textbackslash abs}, \texttt{\textbackslash dd}, \texttt{\textbackslash norm}, \texttt{\textbackslash scpr}), with extra wrapper packages (bm, bbm, dsfont, xspace, cancel, mathtools, cleveref, enumitem). The delivered numbering is the unit's own. Assets: no retained span contains a figure environment, a graphics-inclusion command, a PDF-page inclusion, a table or a TikZ picture; no image file is copied into this package, the package declares no asset dependency and no image file is redistributed with it. Licence: version arXiv:2507.14499v1 of this article is distributed under the Creative Commons Attribution 4.0 International licence, as recorded in the arXiv licence metadata for this exact version and shown on the abstract page https://arxiv.org/abs/2507.14499v1. A full rights notice is delivered at licenses/arxiv-2507.14499-CC-BY-4.0.txt. Characters: every retained excerpt is pure ASCII, so the delivered engine, XeLaTeX with its default Latin Modern text font, reports no missing character.
\begin{assumption}[Regularity Conditions]
\label{ass:wellposedness_multidim}
The driver $f_\theta: [0,T] \times \R^k \times \R^k \times \R^{k \times d} \to \R^k$ and terminal condition $\xi$ satisfy:
\begin{enumerate}[label=(\roman*)]
    \item \textbf{Measurability and Continuity:} $(t,x,y,z) \mapsto f_\theta(t,x,y,z)$ is jointly measurable, and for a.e. $t \in [0,T]$, it is continuous in $(x,y,z)$.
    \item \textbf{Monotonicity in $y$:} There exists a constant $\mu \in \R$ such that for all $(t,x,z)$ and all $y_1, y_2 \in \R^k$:
    \[ \scpr{y_1 - y_2}{f_\theta(t,x,y_1,z) - f_\theta(t,x,y_2,z)} \le \mu \norm{y_1 - y_2}^2. \]
    \item \textbf{Quadratic Growth in $z$:} There exists a function $\kappa: \R^+ \to \R^+$ and a constant $\alpha \ge 0$ such that for all arguments,
    \[ \norm{f_\theta(t,x,y,z)} \le \kappa(\norm{x}) + \kappa(\norm{y}) + \frac{\alpha}{2}\norm{z}_F^2, \]
    where $\norm{\cdot}_F$ is the Frobenius norm on $\R^{k \times d}$.
    \item \textbf{Integrability:} The terminal value $\xi$ is bounded. The exogenous process $X$ has paths in a compact set, i.e., $\sup_{t \in [0,T]} \norm{X_t(\omega)} \le C_X < \infty$ a.s.
\end{enumerate}
\end{assumption}

\begin{theorem}[Existence of Well-Behaved Neural Drivers]
\label{thm:existence_of_theta}
The class of neural drivers that satisfy both the multi-dimensional well-posedness conditions of \Cref{ass:wellposedness_multidim} and the unique implicit volatility condition of \Cref{ass:implicit_vol} (adapted to the multi-dimensional case) is non-empty. Specifically, one can define families of neural network architectures parameterized by $\theta$ such that any choice of $\theta$ within the family produces a driver with the required properties.
\end{theorem}

\begin{assumption}[Existence of a Unique Implicit Volatility Function]
\label{ass:implicit_vol}
Let the specialized driver $g_\theta: [0,T] \times \R \times \R \to \R$ be of class $C^1$. We say that $g_\theta$ admits a \textbf{unique implicit volatility function} if there exists a function $\nu_\theta: [0,T] \times \R \to (0, \infty)$ satisfying for each $(t,x) \in [0,T] \times \R$:
\begin{enumerate}[label=(\alph*)]
    \item \textbf{Unique Root:} $z = \nu_\theta(t,x)$ is the unique positive real solution to the algebraic equation $g_\theta(t,x,z) = 0$.
    \item \textbf{Regularity at the Root:} $\partial_z g_\theta(t,x,\nu_\theta(t,x)) \neq 0$. By the Implicit Function Theorem, this ensures $\nu_\theta(t,x)$ is of class $C^1$ and thus locally Lipschitz in $x$.
    \item \textbf{Global Linear Growth:} There exists a constant $C > 0$ such that $|\nu_\theta(t,x)| \le C(1+|x|)$ for all $(t,x)$.
\end{enumerate}
\end{assumption}

\begin{assumption}[Lipschitz Regularity of the Implicit Volatility]
\label{ass:mckean_lipschitz_revised}
The implicit volatility function $\nu_\theta: [0,T] \times \R \times \Pcal_2(\R) \to \R$ is measurable and satisfies the following conditions for some constants $L_x, L_\mu, C > 0$:
\begin{enumerate}[label=(\roman*)]
    \item \textbf{Lipschitz continuity in state:} For any $t \in [0,T]$, $x, y \in \R$, and $\mu \in \Pcal_2(\R)$:
    \[ \abs{\nu_\theta(t, x, \mu) - \nu_\theta(t, y, \mu)} \le L_x \abs{x - y}. \]
    \item \textbf{Lipschitz continuity in measure:} For any $t \in [0,T]$, $x \in \R$, and $\mu, \pi \in \Pcal_2(\R)$:
    \[ \abs{\nu_\theta(t, x, \mu) - \nu_\theta(t, x, \pi)} \le L_\mu W_2(\mu, \pi). \]
    \item \textbf{Linear growth:} For any $(t, x, \mu) \in [0,T] \times \R \times \Pcal_2(\R)$:
    \[ \abs{\nu_\theta(t, x, \mu)}^2 \le C^2(1 + \abs{x}^2 + W_2(\mu, \delta_0)^2), \]
    where $\delta_0$ is the Dirac measure at the origin.
\end{enumerate}
\end{assumption}

\begin{theorem}[Existence of Well-Behaved Mean-Field Neural Drivers]
\label{thm:mckean_driver_existence}
The class of neural drivers $g_\theta$ for which the implicitly defined volatility function $\nu_\theta(t,x,\mu)$ satisfies the Lipschitz regularity and growth conditions of \Cref{ass:mckean_lipschitz_revised} is non-empty.
\end{theorem}

\begin{proof}
The proof is constructive. We will define a specific architecture for the specialized driver $g_\theta(t,x,\mu,z)$ and show that by imposing architecturally enforceable constraints on its components, the resulting implicit volatility function $\nu_\theta$ satisfies the conditions of \Cref{ass:mckean_lipschitz_revised}.

\paragraph{Step 1: Architectural Design.}
We structure the specialized driver $g_\theta$ to separate the dependencies on $z$ and $\mu$:
\begin{equation}
\label{eq:mckean_driver_arch}
    g_\theta(t, x, \mu, z) \coloneqq h_{\theta_1}(t, x, z) - \mathbb{E}_{Y \sim \mu}[\phi_{\theta_2}(x, Y)],
\end{equation}
where $h_{\theta_1}: [0,T] \times \R \times \R \to \R$ and $\phi_{\theta_2}: \R \times \R \to \R$ are functions parameterized by distinct neural networks with parameters $\theta_1$ and $\theta_2$ respectively, with $\theta = (\theta_1, \theta_2)$. The expectation $\mathbb{E}_{Y \sim \mu}[\cdot]$ is equivalent to the integral $\int_\R \phi_{\theta_2}(x, y)\mu(\dd y)$.

\paragraph{Step 2: Defining the Implicit Volatility Function $\nu_\theta$.}
The canonical NBM condition is $g_\theta(t, x, \mu, z) = 0$. This yields the equation:
\[ h_{\theta_1}(t, x, z) = \mathbb{E}_{Y \sim \mu}[\phi_{\theta_2}(x, Y)]. \]
To ensure a unique, well-behaved root $z = \nu_\theta(t,x,\mu)$, we enforce structural properties on $h_{\theta_1}$. As in the proof of \Cref{thm:existence_of_theta}, we design the network for $h_{\theta_1}$ such that for any fixed $(t,x)$, the function $z \mapsto h_{\theta_1}(t,x,z)$ is strictly monotonic and surjective onto $\R$. This can be achieved, for example, by ensuring $\partial_z h_{\theta_1}$ is bounded below by a positive constant.

Under this condition, $h_{\theta_1}$ has a well-defined inverse with respect to its third argument, which we denote $h_{\theta_1}^{-1}(t,x,\cdot)$. The implicit volatility function is then given explicitly by:
\begin{equation}
\label{eq:nu_theta_explicit}
    \nu_\theta(t, x, \mu) \coloneqq h_{\theta_1}^{-1}\left(t, x, \mathbb{E}_{Y \sim \mu}[\phi_{\theta_2}(x, Y)]\right).
\end{equation}

\paragraph{Step 3: Imposing Enforceable Architectural Constraints.}
We now impose regularity conditions on the component networks $h_{\theta_1}$ and $\phi_{\theta_2}$. These conditions can be enforced during training using techniques like spectral normalization (to control Lipschitz constants) or by specific architectural choices (e.g., parameterizing weights as positive).
\begin{enumerate}[label=(C\arabic*), leftmargin=*]
    \item \textbf{Monotonicity and Regularity of $h_{\theta_1}$:} For all $(t,x,z)$, there exists a constant $c_h > 0$ such that $\partial_z h_{\theta_1}(t,x,z) \ge c_h$. This ensures the inverse $h_{\theta_1}^{-1}$ exists and, by the inverse function theorem, its partial derivative with respect to its third argument is bounded by $1/c_h$.
    \item \textbf{Lipschitzness of $h_{\theta_1}^{-1}$:} The inverse function $v(t,x,w) \mapsto h_{\theta_1}^{-1}(t,x,w)$ is Lipschitz in all its arguments on its domain. Let its Lipschitz constants be $L_{h,t}, L_{h,x}, L_{h,w}$.
    \item \textbf{Lipschitzness of $\phi_{\theta_2}$:} The function $\phi_{\theta_2}(x,y)$ is globally Lipschitz in both its arguments, with constants $L_{\phi,x}$ and $L_{\phi,y}$. That is, for any $x_1, x_2, y_1, y_2 \in \R$:
    \[ |\phi_{\theta_2}(x_1, y_1) - \phi_{\theta_2}(x_2, y_2)| \le L_{\phi,x}|x_1-x_2| + L_{\phi,y}|y_1-y_2|. \]
    \item \textbf{Growth of $\phi_{\theta_2}$:} $\phi_{\theta_2}$ has at most linear growth: there is a constant $C_\phi > 0$ such that $|\phi_{\theta_2}(x,y)| \le C_\phi(1+|x|+|y|)$.
\end{enumerate}

\paragraph{Step 4: Proof of Lipschitz Continuity in Measure (Condition ii).}
Let $(t,x)$ be fixed. We want to bound $|\nu_\theta(t, x, \mu_1) - \nu_\theta(t, x, \mu_2)|$ for $\mu_1, \mu_2 \in \Pcal_2(\R)$.
Let $I(\mu) \coloneqq \mathbb{E}_{Y \sim \mu}[\phi_{\theta_2}(x, Y)]$. From \eqref{eq:nu_theta_explicit}, we have:
\begin{align*}
    |\nu_\theta(t, x, \mu_1) - \nu_\theta(t, x, \mu_2)| &= |h_{\theta_1}^{-1}(t,x,I(\mu_1)) - h_{\theta_1}^{-1}(t,x,I(\mu_2))| \\
    &\le L_{h,w} |I(\mu_1) - I(\mu_2)| \quad \text{(by (C2))}.
\end{align*}
Now we bound $|I(\mu_1) - I(\mu_2)|$. By the definition of the 2-Wasserstein distance, there exists a coupling $(Y_1, Y_2)$ of random variables with marginal laws $\mu_1$ and $\mu_2$ respectively, such that $W_2(\mu_1, \mu_2)^2 = \E[|Y_1-Y_2|^2]$.
\begin{align*}
    |I(\mu_1) - I(\mu_2)| &= |\E_{Y_1 \sim \mu_1}[\phi_{\theta_2}(x, Y_1)] - \E_{Y_2 \sim \mu_2}[\phi_{\theta_2}(x, Y_2)]| \\
    &= |\E[\phi_{\theta_2}(x, Y_1) - \phi_{\theta_2}(x, Y_2)]| \quad \text{(using the coupling)} \\
    &\le \E[|\phi_{\theta_2}(x, Y_1) - \phi_{\theta_2}(x, Y_2)|] \quad \text{(by Jensen's inequality)} \\
    &\le \E[L_{\phi,y}|Y_1 - Y_2|] \quad \text{(by Lipschitzness of $\phi_{\theta_2}$ in $y$, (C3))} \\
    &\le L_{\phi,y} \left( \E[|Y_1 - Y_2|^2] \right)^{1/2} \quad \text{(by Cauchy-Schwarz)} \\
    &= L_{\phi,y} W_2(\mu_1, \mu_2).
\end{align*}
Combining the inequalities, we get:
\[ |\nu_\theta(t, x, \mu_1) - \nu_\theta(t, x, \mu_2)| \le (L_{h,w} L_{\phi,y}) W_2(\mu_1, \mu_2). \]
This establishes Lipschitz continuity in the measure with constant $L_\mu = L_{h,w} L_{\phi,y}$.

\paragraph{Step 5: Proof of Lipschitz Continuity in State (Condition i).}
Let $(t,\mu)$ be fixed. Let $I(x,\mu) \coloneqq \mathbb{E}_{Y \sim \mu}[\phi_{\theta_2}(x, Y)]$.
\begin{align*}
    &|\nu_\theta(t, x_1, \mu) - \nu_\theta(t, x_2, \mu)| = |h_{\theta_1}^{-1}(t,x_1,I(x_1,\mu)) - h_{\theta_1}^{-1}(t,x_2,I(x_2,\mu))| \\
    &\le L_{h,x}|x_1-x_2| + L_{h,w}|I(x_1,\mu) - I(x_2,\mu)| \quad \text{(by (C2))}.
\end{align*}
We bound the second term:
\begin{align*}
    |I(x_1,\mu) - I(x_2,\mu)| &= |\E_{Y \sim \mu}[\phi_{\theta_2}(x_1, Y) - \phi_{\theta_2}(x_2, Y)]| \\
    &\le \E_{Y \sim \mu}[|\phi_{\theta_2}(x_1, Y) - \phi_{\theta_2}(x_2, Y)|] \\
    &\le \E_{Y \sim \mu}[L_{\phi,x}|x_1 - x_2|] = L_{\phi,x}|x_1 - x_2| \quad \text{(by (C3))}.
\end{align*}
Substituting this back, we get:
\[ |\nu_\theta(t, x_1, \mu) - \nu_\theta(t, x_2, \mu)| \le (L_{h,x} + L_{h,w}L_{\phi,x})|x_1-x_2|. \]
This establishes Lipschitz continuity in the state with constant $L_x = L_{h,x} + L_{h,w}L_{\phi,x}$.

\paragraph{Step 6: Proof of Linear Growth (Condition iii).}
We need to bound $|\nu_\theta(t,x,\mu)|^2$. We assume without loss of generality that $h_{\theta_1}^{-1}(t,x,0)=0$ and $\phi_{\theta_2}(0,0)=0$.
\begin{align*}
    |\nu_\theta(t,x,\mu)| &= |h_{\theta_1}^{-1}(t,x,I(x,\mu)) - h_{\theta_1}^{-1}(t,x,0)| \\
    &\le L_{h,x}|x| + L_{h,w}|I(x,\mu)|.
\end{align*}
Now we bound $|I(x,\mu)|$ using the linear growth condition (C4) on $\phi_{\theta_2}$:
\begin{align*}
|I(x,\mu)| &= |\E_{Y \sim \mu}[\phi_{\theta_2}(x,Y)]| \le \E_{Y \sim \mu}[|\phi_{\theta_2}(x,Y)|] \\
&\le \E_{Y \sim \mu}[C_\phi(1+|x|+|Y|)] = C_\phi(1+|x|+\E_{Y \sim \mu}[|Y|]).
\end{align*}
By Cauchy-Schwarz, $\E[|Y|] \le (\E[|Y|^2])^{1/2} = (\int y^2 \mu(\dd y))^{1/2} = W_2(\mu,\delta_0)$.
So, $|I(x,\mu)| \le C_\phi(1+|x|+W_2(\mu,\delta_0))$. Substituting back:
\[ |\nu_\theta(t,x,\mu)| \le L_{h,x}|x| + L_{h,w}C_\phi(1+|x|+W_2(\mu,\delta_0)) \le C(1+|x|+W_2(\mu,\delta_0)), \]
for some constant $C$. Squaring this gives:
\[ |\nu_\theta(t,x,\mu)|^2 \le C^2(1+|x|+W_2(\mu,\delta_0))^2 \le 3C^2(1+|x|^2+W_2(\mu,\delta_0)^2), \]
which is the required linear growth condition.

Since we have constructively shown that there exists a family of neural network architectures for $g_\theta$ such that the resulting implicit volatility function $\nu_\theta$ satisfies all three conditions of \Cref{ass:mckean_lipschitz_revised}, the set of such drivers is non-empty. This completes the proof.
\end{proof}

\end{document}
